\documentclass{beamer}
\usetheme{Madrid}
\usepackage[utf8]{inputenc}
\usepackage{hyperref}
\title{Digital Analysis of Eusebius}
\author{COMP3420 Documentation}
\date{\today}
\begin{document}
\begin{frame}
  \titlepage
\end{frame}

\begin{frame}{Project Overview}
  \begin{itemize}
    \item Eusebius' \emph{Description of Greece} (2nd~c.~CE) mixes myth and travelogue.
    \item Scholars debate how critically Eusebius reports local traditions.
    \item This project provides tools to analyse the text using modern NLP.
    \item Pipeline: import \textrightarrow{} annotate \textrightarrow{} model \textrightarrow{} visualise.
  \end{itemize}
\end{frame}

\begin{frame}{Data Ingestion}
  \begin{itemize}
    \item Raw text segmented as passages with IDs (e.g., \texttt{1.1.1}).
    \item \texttt{eusebius\_importer.py} loads passages into a SQLite database.
    \item \texttt{translate\_eusebius.py} uses the OpenAI API to produce English translations.
  \end{itemize}
\end{frame}

\begin{frame}{Proper Noun Extraction}
  \begin{itemize}
    \item \texttt{extract\_proper\_nouns.py} queries the API with tool calls to identify people, places, and deities.
    \item Results stored with entity type and canonical form.
    \item These nouns act as stopwords and later form a co--occurrence network.
  \end{itemize}
\end{frame}

\begin{frame}{Machine Learning Models}
  \begin{itemize}
    \item Logistic regression models are built using \texttt{scikit\-learn}.
    \item COMP3420 students will recognise this pipeline of vectorisation and classification.
  \end{itemize}
\end{frame}

\begin{frame}{Network Analysis}
  \begin{itemize}
    \item \texttt{analyse\_noun\_network.py} constructs a graph with NetworkX.
    \item Nodes: proper nouns; edges: co--occurrence in passages.
    \item Centrality metrics (degree, betweenness, PageRank) reveal key figures like Apollo or Athens.
    \item Output visualisations use D3.js in an interactive web page.
  \end{itemize}
\end{frame}

\begin{frame}{Website Generation}
  \begin{itemize}
    \item Modular code in \texttt{website/} creates a static site showing analyses.
    \item Passages are highlighted with predictive words, translated text, and lists of proper nouns.
    \item \texttt{cronscript.sh} can run the entire pipeline daily and sync results to a server.
  \end{itemize}
\end{frame}

\begin{frame}{Takeaways}
  \begin{itemize}
    \item Combines classical scholarship with modern NLP methods you know from scikit--learn and Keras.
    \item Encourages reflection on how language shapes interpretation of ancient texts.
    \item Code is modular and can be adapted for other corpora or further analyses.
  \end{itemize}
\end{frame}

\end{document}
